%% v1.1 Part-I lookup compression; complete source proofs: annex/p01.
\chapter{Mathematical Preliminaries}
\label{sec:p01-mathematical-preliminaries}
Classical-set-theoretic substrate; citations use \cite{Jech2003}.  Depth rows
D-P1-005 and D-P1-001--004,006 are homed in Parts~2 and~6.
\section{Sets, relations, and functions}
\label{sec:p01-sets-relations-functions}
\begingroup\small
\begin{longtable}{@{}p{.22\textwidth}p{.72\textwidth}@{}}
\caption{Set, relation, and function register.}\label{tab:p01-set-function-register}\\
\toprule Term & Exact clause \\ \midrule
Subset\index{subset} & \(A\subseteq B\iff\forall x(x\in A\to x\in B)\);
extensionality: \(A=B\iff(A\subseteq B\land B\subseteq A)\). \\
Power set\index{power set} & \(\mathcal P(A)\defeq\{X:X\subseteq A\}\). \\
Class\index{class}\index{proper class} & Metalinguistic collection of sets;
proper iff not a set; its quantification is metalinguistic. \\
Ordered pair\index{ordered pair} &
\(\langle a,b\rangle\defeq\{\{a\},\{a,b\}\}\). \\
Product\index{Cartesian product} &
\(A\times B\defeq\{\langle a,b\rangle:a\in A,b\in B\}\); for
\(n\in\mathbb N\),
\(A^n=\{f:f\text{ is a function }n\to A\}\), and
\(A^0=\{\varnothing\}\). \\
Relation\index{binary relation} & \(R\) is a relation from \(A\) to \(B\)
iff \(R\subseteq A\times B\), and on \(A\) iff \(R\subseteq A^2\). \\
Properties\index{relation!properties} & Reflexive \(\forall a\,aRa\);
irreflexive \(\forall a\,\neg aRa\); symmetric: if \(aRb\), then \(bRa\);
transitive: if \(aRb\) and \(bRc\), then \(aRc\) (variables in \(A\)). \\
Orders\index{strict linear order}\index{well-order} & A strict linear order is
irreflexive and transitive, with exactly one direction between distinct
elements.  A strict well-order is such an order in which each nonempty subset
has an \(R\)-least member. \\
Function\index{function}\index{range of a function} & \(f:A\to B\):
\(f\subseteq A\times B\), and for every \(a\in A\) there is exactly one
\(b\in B\) such that \(\langle a,b\rangle\in f\); range
\(\{f(a):a\in A\}\) need not equal \(B\).
\(f:A\rightharpoonup B\) means \(f:D\to B\), stated \(D\subseteq A\).
\index{partial function} \\
Map properties\index{injection}\index{surjection}\index{bijection} &
For \(a,a'\in A\), injective: if \(f(a)=f(a')\), then \(a=a'\); surjective:
\(\forall b\in B\,\exists a\in A\ f(a)=b\); bijective: both. \\
Image\index{image}\index{inverse image} & For \(X\subseteq A,Y\subseteq B\),
\(f[X]=\{f(x):x\in X\}\) and
\(f^{-1}[Y]=\{x\in A:f(x)\in Y\}\); the latter need not denote an inverse. \\
Equivalence\index{equivalence relation} & Reflexive, symmetric, transitive
\(E\subseteq A^2\). \\
Quotient\index{equivalence class}\index{quotient set} & For an equivalence
relation \(E\) on \(A\), \([a]_E=\{x\in A:xEa\}\) and
\(A/E=\{[a]_E:a\in A\}\). \\
\bottomrule
\end{longtable}
\endgroup

\begin{hbwarning}{def:p01-set-builder-existence}{Set-builder existence}
\(\{x:P(x)\}\) may be proper; supply a bound or set-existence axiom.
\end{hbwarning}
\begin{hbmethod}{def:p01-quotient-definition}{Definition on a quotient}
\(\bar f([a]_E)=f(a)\) defines \(\bar f:A/E\to B\) only if
\(f(a)=f(a')\) whenever \(aEa'\).
\end{hbmethod}
\begin{hblemma}{thm:p01-ordered-pair-identity}{proved=annex:p1.1}{Ordered-pair identity}
For all sets \(a,b,c,d\),
\[
\langle a,b\rangle=\langle c,d\rangle
\iff(a=c\land b=d).
\]
\end{hblemma}
\section{Size, ordinals, and cardinal comparison}
\label{sec:p01-size-ordinals}
\begin{hbconvention}{def:p01-natural-numbers}{Natural numbers and finite ordinals}
\(0\in\mathbb N\); for \(n\in\mathbb N\),
\(n=\{0,\ldots,n-1\}\); a natural is positive iff \(>0\).
A set is finite iff bijective with some \(n\in\mathbb N\), countable iff it
injects into \(\mathbb N\), and countably infinite iff bijective with
\(\mathbb N\).
\index{natural number}\index{finite ordinal}\index{positive integer}
\index{finite set}\index{countable set}\index{countably infinite set}
\end{hbconvention}
\begingroup\small
\begin{longtable}{@{}p{.22\textwidth}p{.72\textwidth}@{}}
\caption{Cardinality and ordinal register.}\label{tab:p01-cardinal-ordinal-register}\\
\toprule Term & Exact clause \\ \midrule
Comparison\index{cardinality!comparison} & \(\abs A\leq\abs B\) iff an
injection \(A\to B\) exists; \(\abs A=\abs B\) iff a bijection exists;
\(\abs A<\abs B\) iff the first holds and equality does not.
\(\abs{\cdot}\) means cardinality only. \\
Transitive\index{transitive set} &
\(\forall x\in A\,\forall y\in x\ y\in A\), equivalently
\(\bigcup A\subseteq A\). \\
Ordinal\index{ordinal} & Transitive \(\alpha\), strictly \(\in\)-well-ordered;
\(\alpha<\beta\iff\alpha\in\beta\). \\
Cardinal\index{cardinal}\index{infinite cardinal} & Ordinal \(\kappa\) not
bijective with a smaller one; \(\kappa\) is infinite iff it is not bijective
with any \(n\in\mathbb N\). \\
Successor; limit\index{successor ordinal}\index{limit ordinal} &
\(\alpha+1=\alpha\cup\{\alpha\}\); a limit is a nonzero nonsuccessor. \\
\bottomrule
\end{longtable}
\endgroup

\begin{hbtheorem}{thm:p01-csb}{proved=annex:p1.2}{Cantor--Schr\"oder--Bernstein theorem}
For sets \(A,B\), injections \(A\to B\) and \(B\to A\) imply a bijection
\(A\to B\).
\end{hbtheorem}
\begin{hbtheorem}{thm:p01-cantor}{proved=annex:p1.3}{Cantor's theorem}
For every set \(A\), no surjection \(A\to\mathcal P(A)\) exists; hence
\(\abs A<\abs{\mathcal P(A)}\).
\end{hbtheorem}
\begin{hbtheorem}{thm:p01-transfinite-recursion}{cited=Jech2003}{Transfinite recursion}
Let \(\delta\) be ordinal and functional \(G\) assign a set \(G(g)\) to every
function \(g\) with ordinal domain below \(\delta\).  A unique function
\(F\), domain \(\delta\), satisfies
\(F(\alpha)=G(F\mathbin{\upharpoonright}\alpha)\) for all \(\alpha<\delta\).
This is a metalinguistic principle, not an object-language rule.
\end{hbtheorem}

\section{Choice principles and finite coding}
\label{sec:p01-choice-coding}
\begin{hbdefinition}{def:p01-axiom-choice}{Axiom of Choice}
Every set \(X\) of nonempty sets has \(c:X\to\bigcup X\) with
\(c(A)\in A\) for all \(A\in X\).
\end{hbdefinition}
\begin{hbtheorem}{thm:p01-ultrafilter-lemma}{cited=HalpernLevy1971}{Ultrafilter Lemma}
Let \(I\) be a set.  Every \(\mathcal F\subseteq\mathcal P(I)\) containing
\(I\), omitting \(\varnothing\), and upward- and finite-intersection-closed
extends to \(\mathcal U\subseteq\mathcal P(I)\) with those properties and
exactly one of \(X,I\setminus X\) for each \(X\subseteq I\).
This metatheoretic choice principle follows from Choice and, over
\(\mathsf{ZF}\), is equivalent to the Boolean Prime Ideal Theorem
\cite{Jech2003}.  Halpern and L\'evy prove that if \(\mathsf{ZF}\) is
consistent, then so is \(\mathsf{ZF}\) together with this lemma and the
failure of Choice \cite{HalpernLevy1971}; thus, on that consistency
hypothesis, the lemma does not imply Choice.  This is a relative-consistency
claim, not absolute independence.  D-P1-005 gives the Part~2 compactness link.
\end{hbtheorem}
\begin{hbtheorem}{thm:p01-well-ordering}{cited=Jech2003}{Well-ordering theorem}
Choice implies that every set can be well-ordered.
\end{hbtheorem}
\begin{hbcorollary}{thm:p01-cardinal-assignment}{cited=Jech2003}{Cardinal assignment}
Choice implies every set \(A\) is bijective with a unique cardinal
\(\abs A\); \(\abs{\mathbb N}=\aleph_0\).
\end{hbcorollary}
\begin{hbconvention}{def:p01-cardinal-maximum}{Maximum of a finite cardinal list}
For a nonempty finite list of cardinals \(\kappa_1,\ldots,\kappa_n\),
\(\max(\kappa_1,\ldots,\kappa_n)\) is its greatest member.  Choice makes
comparison total.  Henkin and
L\"owenheim--Skolem constructions use this convention.
\index{maximum!finite cardinal list}
\end{hbconvention}
For any set \(A\),
\(A^{<\omega}\defeq\bigcup_{n\in\mathbb N}A^n\) is its set of finite
sequences.  When \(A\) is an alphabet, the string notation \(A^*\) below
denotes the same set.
\begin{hbtheorem}{thm:p01-finite-coding}{cited=Jech2003}{Finite coding over an infinite cardinal}
Assume Choice and let \(\kappa\) be an infinite cardinal.  For positive
\(n\in\mathbb N\),
\(\abs{\kappa^n}=\kappa\) and \(\abs{\kappa^{<\omega}}=\kappa\); if
\(\alpha<\kappa\) and each \(A_\beta\) (\(\beta<\alpha\)) is finite, then
\(\abs{\bigcup_{\beta<\alpha}A_\beta}<\kappa\); if
\(\abs{A_n}\leq\kappa\) for every \(n\in\mathbb N\), then
\(\abs{\bigcup_{n\in\mathbb N}A_n}\leq\kappa\).  This excludes arbitrary
infinite products.
\end{hbtheorem}
\begin{hbcorollary}{thm:p01-syntax-counting}{proved=annex:p1.4}{Syntax counting}
Assume Choice, let \(\kappa\) be an infinite cardinal, and suppose
\(\abs A\leq\kappa\).  Then
\(\abs{A^{<\omega}}\leq\kappa\).  Alphabets of that size, their finite
strings and tuples, and expressions of fixed finitary grammars adding at most
countably many symbols have size at most \(\kappa\).
\end{hbcorollary}

\section{Finite strings, induction, and closure}
\label{sec:p01-strings-recursion}
\begingroup\small
\begin{longtable}{@{}p{.22\textwidth}p{.72\textwidth}@{}}
\caption{Finite-string and closure register.}\label{tab:p01-string-register}\\
\toprule Term & Exact clause \\ \midrule
Sequence\index{finite sequence} & For a set \(A\) and \(n\in\mathbb N\), a
sequence in \(A\) of length \(n\) is a function \(n\to A\);
\(\epsilon\) is the unique length-zero sequence. \\
String\index{string} & Finite sequence over \(A\);
\(A^*=\bigcup_{n\in\mathbb N}A^n\).  \(*\) requires string scope. \\
Concatenation\index{concatenation} & For \(u\in A^m,v\in A^n\):
\((uv)(i)=u(i)\) if \(i<m\), and \(v(i-m)\) if \(m\leq i<m+n\);
\(uv\in A^{m+n}\). \\
Generated set\index{inductive definition} & Given set \(U\), \(B\subseteq U\),
and set \(F\) of finitary operations on \(U\):
\(\bigcap\{X\subseteq U:B\subseteq X,\ X\text{ is closed under every
operation in }F\}\). \\
Reflexive-transitive closure\index{reflexive transitive closure} & Least
reflexive transitive \(R^*\supseteq R\) on \(X\), for \(R\subseteq X^2\). \\
Equivalence closure\index{equivalence closure} & Least equivalence
\(E\supseteq R\) on \(X\), for \(R\subseteq X^2\). \\
\bottomrule
\end{longtable}
\endgroup

The countability proof below requires a map coding two natural numbers by one;
the set-valued Kuratowski ordered pair cannot serve because its codomain is not
\(\mathbb N\).  As a metalinguistic numerical-coding abbreviation, define
\[
 \operatorname{pair}_{\mathbb N}:\mathbb N^2\to\mathbb N,
 \qquad
 \operatorname{pair}_{\mathbb N}(a,b)
 \defeq\frac{(a+b)(a+b+1)}2+b.
\]
\index{pairing function!natural numbers}
Annex~P1.5 proves bijectivity before use; this is neither the Kuratowski pair
nor a sequence notion.
\begin{hblemma}{thm:p01-countable-strings}{proved=annex:p1.5}{Countability of finite strings}
If \(A\) is countable, then \(A^*\) is countable.
\end{hblemma}
\begin{hbcorollary}{thm:p01-countable-formulas}{proved=annex:p1.6}{Countability of formulas over a countable alphabet}
Formulas generated as strings over a countable alphabet form a countable set.
\end{hbcorollary}
\begin{hblemma}{thm:p01-concatenation-laws}{proved=annex:p1.7}{Concatenation laws}
For \(u,v,w\in A^*\): \((uv)w=u(vw)\), \(\epsilon u=u\), \(u\epsilon=u\).
\end{hblemma}
\begin{hbtheorem}{thm:p01-structural-induction}{proved=annex:p1.8}{Structural induction}
If \(G\) is generated from bases \(B\) by finitary constructors \(F\), and
the metalinguistic property \(P\) of members of \(G\) holds on \(B\) and is
preserved by each constructor from all immediate arguments, then
\(\forall x\in G\ P(x)\).
\end{hbtheorem}
\begin{hbtheorem}{thm:p01-structural-recursion}{proved=annex:p1.9}{Structural recursion}
Let \(G\) be inductively generated from the declared bases and finitary
constructors, and suppose each member has a unique construction.  A value in
\(Y\) for each base and a map \(H_f:Y^k\to Y\) for each \(k\)-ary constructor
determine a unique \(h:G\to Y\) with
\[
h(f(x_1,\ldots,x_k))=H_f(h(x_1),\ldots,h(x_k)).
\]
\end{hbtheorem}
\begin{hbmethod}{def:p01-derivation-induction}{Induction on a derivation}
Use the inductive definition or state a measure; check each initial derivation
and rule from its immediate subderivations, including freshness, discharge,
and other side conditions.  For soundness, prove local semantic preservation.
\end{hbmethod}
\begin{hbwarning}{def:p01-structural-induction-cases}{The active grammar fixes the cases}
Check every active base and constructor; another grammar's cases do not suffice.
\end{hbwarning}
\begin{hblemma}{thm:p01-rtc-path}{proved=annex:p1.10}{Reflexive-transitive closure paths}
For \(R\subseteq X^2\) and \(x,y\in X\), \(xR^*y\) iff some
\(n\in\mathbb N\) and \(x_0,\ldots,x_n\in X\) satisfy \(x_0=x\), \(x_n=y\),
and \(x_iRx_{i+1}\) for every \(i<n\).
\end{hblemma}
